% Common preamble for Erdos 304 Round 4 (XeLaTeX)
\documentclass[11pt,a4paper]{article}
\usepackage{fontspec}
\setmainfont{Noto Serif CJK KR}[ItalicFont={Noto Serif CJK KR},ItalicFeatures={FakeSlant=0.18},BoldItalicFont={Noto Serif CJK KR Bold},BoldItalicFeatures={FakeSlant=0.18}]
\setsansfont{Noto Sans CJK KR}
\XeTeXlinebreaklocale "ko"
\XeTeXlinebreakskip = 0pt plus 1pt
\usepackage{newunicodechar}
\newfontfamily\latinfb{DejaVu Serif}[Scale=MatchLowercase]
\newunicodechar{ő}{{\latinfb\symbol{"0151}}}
\newunicodechar{é}{{\latinfb\symbol{"00E9}}}
\newunicodechar{è}{{\latinfb\symbol{"00E8}}}
\newunicodechar{ć}{{\latinfb\symbol{"0107}}}
\newunicodechar{č}{{\latinfb\symbol{"010D}}}
\newunicodechar{š}{{\latinfb\symbol{"0161}}}
\newunicodechar{ö}{{\latinfb\symbol{"00F6}}}
\newunicodechar{ü}{{\latinfb\symbol{"00FC}}}
\newunicodechar{á}{{\latinfb\symbol{"00E1}}}
\usepackage{amsmath,amssymb,amsthm,mathtools}
\usepackage[margin=2.3cm]{geometry}
\usepackage{xcolor}
\usepackage{booktabs,longtable,array}
\usepackage{enumitem}
\usepackage[colorlinks=true,linkcolor=blue!50!black,citecolor=green!40!black,urlcolor=blue!60!black]{hyperref}
\usepackage{tcolorbox}
% status tags
\newcommand{\PROVED}{\textcolor{green!45!black}{\textbf{[PROVED]}}}
\newcommand{\COND}[1]{\textcolor{orange!80!black}{\textbf{[CONDITIONAL on #1]}}}
\newcommand{\OPEN}{\textcolor{red!70!black}{\textbf{[OPEN]}}}
\newcommand{\REFUTED}{\textcolor{purple!80!black}{\textbf{[REFUTED]}}}
\newcommand{\NUMERIC}{\textcolor{blue!70!black}{\textbf{[NUMERICAL]}}}
\newcommand{\HEUR}{\textcolor{gray!70!black}{\textbf{[HEURISTIC]}}}
\newcommand{\EXT}{\textcolor{teal!70!black}{\textbf{[EXTERNAL INPUT]}}}
% theorem environments (Korean names)
\theoremstyle{plain}
\newtheorem{theorem}{정리}[section]
\newtheorem{lemma}[theorem]{보조정리}
\newtheorem{proposition}[theorem]{명제}
\newtheorem{corollary}[theorem]{따름정리}
\newtheorem{conjecture}[theorem]{추측}
\newtheorem{hypothesis}[theorem]{가설}
\theoremstyle{definition}
\newtheorem{definition}[theorem]{정의}
\newtheorem{example}[theorem]{예}
\newtheorem{counterexample}[theorem]{반례}
\theoremstyle{remark}
\newtheorem{remark}[theorem]{주의}
% notation
\newcommand{\N}{\mathbb{N}}
\newcommand{\Z}{\mathbb{Z}}
\newcommand{\R}{\mathbb{R}}
\newcommand{\C}{\mathbb{C}}
\newcommand{\T}{\mathbb{T}}
\newcommand{\Lam}{\Lambda}
\newcommand{\eps}{\varepsilon}
\newcommand{\e}{\mathrm{e}}
\DeclareMathOperator{\lcm}{lcm}
